\begin{center}
    {\large \textbf{Giorno 2} \textbar\ Arrivo al Tetebatu \textbar\ 10 Agosto 2024 \textbar\ \textbf{Tetebatu}, Lombok}
\end{center}

\noindent\rule{\textwidth}{0.4pt}
\begin{figure}[h!] % Portrait images below
    \centering
    % First portrait image (on the left)
    \begin{minipage}[h]{0.48\textwidth} % Set width equal to half the page
        \centering
        \includegraphics[width=\textwidth]{images/flags/Screenshot Tetebatu.png}
    \end{minipage}
    \hfill
    % Text
    \begin{minipage}[h]{0.48\textwidth}
        \raggedright
        \textbf{Superficie}: 4.567 km² \\
        \vspace{0.3cm}
        \textbf{Popolazione}: 3,3 milioni \\
        \vspace{0.3cm}
        \textit{L'isola di Lombok fa parte dell'arcipelago delle Piccole Isole della Sonda. È nota per le spiagge e il surf, specialmente a Kuta e Banko Banko. Le isole Gili, a ovest della costa di Lombok, offrono spiagge, barriere coralline perfette per immersioni e snorkeling, e anche un vivaio per le tartarughe marine.}
    \end{minipage}
\end{figure}


\landscapeLargeMargin{images/day2/PXL_20240810_053956770.jpg}

\landscapeLargeMargin{images/day2/PXL_20240810_083223674.jpg}

\landscapeLargeMargin{images/day2/PXL_20240810_083216835.jpg}

\landscapeLargeMargin{images/day2/PXL_20240810_090659642.jpg}

\landscapeLargeMargin{images/day2/PXL_20240810_090723704.jpg}

\landscapeLargeMargin{images/day2/PXL_20240810_092451778.MP.jpg}

\landscapeLargeMargin{images/day2/PXL_20240810_100740449.NIGHT.jpg}

\noindent 10 Agosto 2024 \newline

Dopo l'infinita giornata di ieri, quel che ci vuole è proprio un bel sonno ristoratore. All'una del mattino sono già sveglio a causa del fuso orario completamente sballato. Sento la stanchezza del corpo, le membra pesanti e quel malessere allo stomaco che mi ricorda quando facevo le notti in bianco prima degli esami ingurgitando solo caffè, ma l'unica cosa che non sento è il sonno. Mi giro, mi rigiro, penso a come stia Ramon, mi giro e mi rilasso, mi bagno, m'asciugo e inizia qui lo spasso — solo che di spasso proprio non se ne parla e anzi sale solo l'astio.

Mi addormento quando sta per albeggiare, così quando la sveglia suona poco prima delle sette la interrompo per concedermi un minutino di sonno in più. Quando apro gli occhi è passata mezz'ora e l'autista del Grab mi ha già scritto che è quasi arrivato davanti all'hotel. Si tenga conto che la sveglia originale era già stata impostata in modalità Fantozzi: a forza di esperimenti e perfezionamenti continui, l'avevamo messa alle 6:50, vale a dire al limite delle possibilità umane. Ora abbiamo dieci minuti per scendere dal letto, lavarci, vestirci, chiudere i bagagli, lasciare la stanza, scendere nella hall e fare il check-out. Rinunciando a colazione e caffè riusciamo incredibilmente ad arrivare puntuali.

Saliti in auto vorrei solo che l'autista si andasse a schiantare contro un muro per porre fine alle mie sofferenze, ma conoscendomi rimarrei illeso solo per il gusto sadico di non potermi riposare. Il Messico era terminato con la privazione del sonno e questo viaggio inizia già così.

Per fortuna il traffico la mattina è scorrevole e arriviamo subito in aeroporto, dove imbarchiamo i bagagli — dedicando un pensiero al poveretto a cui abbiamo rubato lo zaino la sera prima — e ci prendiamo un po' di tempo per fare colazione in un bar in stile occidentale. Muffin e cappuccino per me, pain au chocolat e tè Jasmine per Elisa.

Con calma ci avviciniamo ai gate, ma non prima di aver fatto tappa da Ralph Lauren: non avrei mai pensato di comprare questi marchi qui in Indonesia, ma una Ralph originale a 18€ fa ridere, dai. Con nonchalance superiamo il gate e ci prepariamo ad aspettare l'imbarco.

Subito dopo di noi arriva il gruppo con sciarpe e borse bianche e rosse che era in volo con noi il giorno prima da Jeddah. Non vedendo la vecchietta realizzo che probabilmente non ce l'ha fatta. Mentre Eli è al bagno, uno del gruppo mi si avvicina — un signore anziano alto quanto un cespuglio — e mi picchietta la gamba ordinando di togliere lo zaino per farlo sedere. Gli dico di no, ma non faccio in tempo a spiegargli che non sono solo che già si indispone e insiste, per poi andarsene rabbioso quando gli spiego che mia moglie sta tornando. Strani incontri.

Chiamano per il nostro volo e siamo i primi a imbarcarci. Non ero mai entrato in un aereo vuoto! Io ed Eli non siamo seduti vicini e approfitto del breve viaggio per finire di leggere un libro e scrivere un po' di diario, ma vengo spesso distratto dalle porno-hostess della compagnia aerea. Penso che a fine servizio dovranno scucirgli la divisa aderente o soffocheranno. Nel frattempo Elisa è in buona compagnia: il passeggero seduto accanto a lei ha un singhiozzo talmente forte che sembra stia per vomitare da un momento all'altro.

Atterriamo a Lombok e dopo una breve sosta a rifocillarci con quello che troviamo nel primo bar disponibile raggiungiamo il nostro autista, che deve portarci ad Al Sasaki Homestay. Facendo due chiacchiere ci spiega che abita a pochi metri da Al Sasaki e che conosce i gestori da quando era piccolo. "They're like family."

Il tragitto dura circa un'ora e mezza e le strade che percorriamo ricordano tanto il Messico, lo Yucatán in particolare, con interi villaggi costruiti lungo il bordo della strada principale. Qui l'asfalto è ancora più dissestato e la guida più folle. Il nostro autista usa costantemente il clacson per avvisare di un sorpasso e le quattro frecce per ringraziare una volta superato. L'unico dettaglio che fa capire che non ci troviamo in Messico sono le moschee e la continua musica da preghiera che sentiamo in sottofondo attraversando i centri abitati più grandi. Altro dettaglio distintivo sono le risaie: l'Indonesia ne è piena e durante il viaggio capita di vedere donne che essiccano il riso su teloni a lato della strada.

A un certo punto il nostro autista fa una deviazione per un motivo che non capisco. O c'è qualche celebrazione in corso che blocca la strada, o — molto più probabilmente — ci sono dei controlli e lui non ha la patente. Mi spiega comunque che è meglio muoversi in macchina la mattina, perché dopo l'una ci sono tutte le cerimonie locali e il traffico aumenta. Eli trascorre metà viaggio addormentata per risvegliarsi esattamente tre minuti prima di arrivare a destinazione. Ormai ha dormito su praticamente ogni mezzo di trasporto da quando siamo partiti.

Scesi dall'auto ci addentriamo nel villaggio da un ingresso a lato della strada e dopo aver percorso — con un misto di eccitazione e trepidazione — strette vie passando davanti a case modeste, ci ritroviamo in una location mozzafiato. Siamo accolti in uno spazio all'aperto con cucina, tavolini, divanetti e sedie di tela sospese, circondato dalla giungla. Ci sono varie villette private per gli ospiti e la nostra è l'unica al primo piano, sopra la cucina, e gode di una vista spettacolare sulla foresta di palme e banani. Ha una porta in legno interamente intagliata con splendide venature, un letto matrimoniale in legno con spalliera decorata come il portone, una doccia realizzata con sassi e legno e un lavandino in pietra.

Appena arrivati ci accoglie il proprietario Ronnie, insieme alle donne di casa — tutte rigorosamente col velo — e Oshi, del quale ancora non abbiamo capito se sia il figlio di Ronnie o solo un fedele dipendente. Ci offrono un cocktail di benvenuto al dragon fruit e banana e ci serviamo di crackers di riso artigianali.

Ci diamo poi una veloce rinfrescata: siamo stanchissimi ma non vogliamo buttare la giornata, così chiediamo subito di poter fare uno dei tour organizzati per le poche ore di luce che restano nel pomeriggio. Anche se non erano entusiasti di doverci accompagnare a fare un tour non preventivato non lo danno a vedere e sono subito disponibili, nonostante la loro proposta iniziale per la giornata fosse stata: "Chill and relax."

È così che veniamo caricati su due motorini che ci portano in giro per tutta la regione del Tetebatu. Elisa con Oshi, io con un vicino che è la metà di me per peso e altezza. Quando sono seduto dietro di lui sembriamo padre e figlio a lezione di guida: potrei appoggiargli il mento sulla nuca da quanto è minuto, ma d'altro canto godo di una visuale piena. Riesco a fare qualche video dello splendido panorama che incontriamo.

Passiamo accanto a risaie, a campi di tabacco, a un campo da calcio dove è in corso una partita tra squadre locali con tifo sfegatato, a villaggi dove tutti i ragazzini ci salutano, a negozi improvvisati di ogni tipo di cibarie e a venditori di bottiglie di benzina. Quest'ultima è un'usanza molto diffusa: in tantissimi vendono benzina e gasolio in bottiglie di vetro a bordo strada, travasata con un imbuto nei serbatoi dei veicoli, soprattutto motorini. Meno di un euro a bottiglia: se al ritorno avrò ancora spazio nel bagaglio a mano potrei farci un pensierino.

La prima tappa è in un villaggio vicino dove tutte le donne si dedicano alla tessitura dei Sarong. Visitiamo prima il negozio espositivo, poi il proprietario ci accompagna a vedere come le donne locali producono i tessuti a mano con gli antichi strumenti di legno. Per un Sarong singolo — largo 60 cm e lungo 4 m — possono impiegare fino a un mese. Il filo è tinto con coloranti artificiali o naturali: questi ultimi, meno luminosi, si ottengono immergendo e lasciando essiccare a più riprese i filamenti in tinte naturali. Con l'ananas si ottiene il giallo, col frutto della passione il rosa, con la corteccia d'albero il marrone, con la pianta dell'indaco il blu. Per il verde si mescolano ananas e indaco.

Terminata la visita risaliamo sui motorini per la seconda tappa, dedicata alla lavorazione del bambù. Una famiglia ci accoglie nella loro dimora all'aperto: ci sediamo tutti scalzi su tappeti di canapa e mentre sorseggiamo il tè offerto da loro e mangiamo crackers di riso e tapioca ci prendono le misure per prepararci degli anelli di bambù intrecciato. Conosciamo il capo famiglia, Mr. Bean, che parla un po' l'italiano grazie a Duolingo e ci insegna parole indonesiane. Impariamo a dire grazie — Terima Kasih — e ovviamente lo dimentichiamo subito. Oltre agli anelli ci regalano un cofanetto di bambù realizzato in pochi minuti davanti a noi, e alla fine compriamo dei braccialetti fatti a mano come ricordo.

Il sole sta ormai tramontando quando li salutiamo e saliamo sui motorini per tornare ad Al Sasaki. Il panorama al tramonto lascia senza fiato. Le discese dell'andata sono diventate salite al ritorno e sento che col mio peso il motorino fatica a salire, mentre il mio autista va a tavoletta e ogni tanto deve sterzare per non farci cadere.

Terminato il tour ringraziamo i nostri driver e ci prepariamo alla cena. Qui funziona così: c'è un menù alla carta con soli piatti locali, si sceglie cosa mangiare e lo si prepara insieme. Elisa è esterrefatta e non vede l'ora di cominciare a cucinare. Scegliamo entrambi i fried noodles — o Mie Goreng — e veniamo stupiti nuovamente dalla meravigliosa semplicità con cui ci fanno preparare questo piatto.

Entriamo in una cucina spartana con pareti e pavimento di pietra, due fornelli da campeggio e un lavandino, dove quattro donne col velo e a piedi nudi sono intente a cucinare e lavare i piatti a mano. Una a una tagliamo tutte le verdure con estrema precisione: sono loro a mostrarci che taglio effettuare e ci correggono strada facendo. Prepariamo spinaci freschi, cavolo, fagiolini e carote, poi mettiamo tutto a friggere in padella e mentre rigiriamo aggiungiamo la polvere e la salsa dei noodles. Nel frattempo una delle donne ha cotto i noodles e li rovescia in padella perché si amalgami tutto. Aggiungiamo la loro salsa di soia, che a differenza della nostra è dolce perché mescolata con zucchero di canna, e impiattiamo aggiungendo pollo, tofu fritto, cetriolo e gallette di semi che nel frattempo le altre donne hanno cucinato. Come contorno, un piatto di patatine fritte.

Terminato di cucinare portiamo i piatti in tavola e iniziamo a gustare la cena. Siamo stanchi e affamati, è vero, ma posso dire con sicurezza che si tratta di uno dei piatti più buoni che abbiamo mai assaggiato. Vorrei godermelo con calma ma devo tenere d'occhio Elisa, che non appena finisce i suoi noodles tenta di rubare i miei. Mentre un pipistrello svolazza sopra le nostre teste ammiriamo la foresta di banani, e forse proprio questo ci ispira: chiediamo come dessert due banane dolcissime e gommose. Sono le banane indonesiane, la specialità Pisang Kepok, più corte e squadrate delle nostre.

Trascorriamo ancora un po' di tempo con il nostro host a parlare del più e del meno. La storia del cavolo va menzionata: gli spieghiamo che vorremmo riprodurre il piatto assaggiato anche a casa, ma il cavolo che ci hanno fatto usare è introvabile in Italia. Ci spiega che a Lombok crescono solo due tipi di cavolo — quello verde, tipico delle zone più calde, e quello rosso, che ha bisogno di temperature più basse e viene coltivato solo in montagna. Questo fa sì che le popolazioni di queste aree, quando i raccolti sono abbondanti, barattino i frutti del proprio lavoro tra di loro, così che tutti possano godere dei sapori della terra. Siamo sempre più affascinati da questa cultura.

Ci salutiamo e decidiamo di fare una passeggiata per le vie del villaggio. Incrociamo pochi abitanti con cui scambiamo sorrisi e saluti, e passando davanti a una casa sentiamo provenire della musica: ci avviciniamo e scopriamo il capo famiglia che suona la chitarra e canta insieme a moglie e figli canzoni dei Beatles. Con questa immagine indelebile torniamo nella giungla e andiamo a provare le amache in fondo al villaggio, da dove ammiriamo il cielo stellato. Qui incontriamo Ronnie, il proprietario, che ci lascia dei dolcetti confezionati al cioccolato. Eli dice di non avere più fame ma ne mangia tre senza che quasi io me ne accorga.

Scattiamo un po' di foto alle stelle e torniamo in camera, pronti finalmente a riposarci. L'indomani ci aspetta un lungo tour per le risaie e le cascate e vogliamo arrivare carichi. Inizio a scrivere questa pagina di diario, ma collasso giusto dopo aver rivissuto l'atterraggio a Lombok.

\pagestyle{empty} % disable numbers

%single
\portraitFullscreen{images/day2/PXL_20240810_082901676.jpg}
%

%coppia
\twoImagesLargeMargins{images/day2/PXL_20240810_090659642.jpg}
                      {images/day2/PXL_20240810_090723704.jpg}
\threeImagesLargeMargins{images/day2/PXL_20240810_100151175.jpg}
                        {images/day2/PXL_20240810_095929164.PORTRAIT.jpg}
                        {images/day2/PXL_20240810_095939186.PORTRAIT.jpg} 
%

%coppia
\landscapeThinMargin{images/day2/PXL_20240810_120038922.MP.jpg}

\twoImagesLargeMargins{images/day2/PXL_20240810_113655543.jpg}
                      {images/day2/PXL_20240810_115606001.MP.jpg}
%

%coppia
\landscapeThinMargin{images/day2/PXL_20240810_114135715.jpg}

\twoImagesLargeMargins{images/day2/PXL_20240810_115853312.MP.jpg} 
                      {images/day2/PXL_20240810_115927640.MP.jpg}
%

%coppia
\twoImagesLargeMargins{images/day2/PXL_20240810_120608464.MP.jpg}
                      {images/day2/PXL_20240810_120615048.jpg}
                      
\twoImagesLargeMargins{images/day2/PXL_20240810_124458770.jpg}
                      {images/day2/PXL_20240810_124542063.jpg} 
%

%coppia
\splitImageLeft{images/day2/PXL_20240810_113053476.jpg}
\splitImageRight{images/day2/PXL_20240810_113053476.jpg}
%

%
\landscapeThinMargin{images/day2/PXL_20240810_131109620.jpg}
%

\pagestyle{fancy} % enable numbers